\documentclass[11pt]{article}
\usepackage[T1]{fontenc}
\usepackage{lmodern,amsmath,amssymb,amsthm,mathtools}
\usepackage[margin=0.85in]{geometry}
\usepackage{booktabs,longtable,array,enumitem,microtype,xcolor,fancyhdr}
\usepackage[hyphens]{url}
\usepackage[colorlinks=true,linkcolor=blue!45!black,urlcolor=blue!45!black,citecolor=blue!45!black]{hyperref}
\hypersetup{pdftitle={Smaller explicit Diophantine forms of ZFC consistency},
 pdfauthor={AI-assisted research contribution for ProveIt},
 pdfsubject={Explicit polynomial constructions and proofs of equivalence to Con(ZFC)},
 pdfkeywords={Diophantine equations, ZFC consistency, Pell equations, arithmetic circuits}}
\usepackage{listings}
\lstset{basicstyle=\ttfamily\small,breaklines=true,columns=fullflexible,keepspaces=true}
\newtheorem{theorem}{Theorem}[section]
\newtheorem{lemma}[theorem]{Lemma}
\newtheorem{proposition}[theorem]{Proposition}
\theoremstyle{remark}\newtheorem{remark}[theorem]{Remark}
\newcommand{\Z}{\mathbb Z}\newcommand{\N}{\mathbb N}
\newcommand{\Pos}{\mathbb Z_{>0}}\newcommand{\Con}{\operatorname{Con}}
\newcommand{\pc}{\operatorname{pc}}\newcommand{\Norm}{\operatorname{N}}
\newcommand{\pin}{\texttt{ae3007ca4467011e0e051c84acb4a1dd13971b0e}}
\setlength{\headheight}{14pt}
\pagestyle{fancy}\fancyhf{}\lhead{Smaller Diophantine forms of ZFC consistency}\rhead{\thepage}
\setlist{itemsep=3pt,topsep=5pt}
\emergencystretch=2em
\title{\textbf{Smaller explicit Diophantine forms\\of ZFC consistency}\\[5pt]
\large Compact root bounds, signed coordinates, and odd scales}
\author{An AI-assisted research contribution prepared for ProveIt}
\date{22 October 2026}
\begin{document}
\maketitle
\begin{abstract}
We combine the current ProveIt constructions with three additional reductions.
The main moderate-degree result is a fully specified polynomial in twelve
positive unknowns of exact degree $1097$, evaluated by $116$ binary arithmetic
operations. Its three-equation presentation has maximum degree $1088$ and costs
$110$ operations. This improves the reviewed twelve-variable degree minimum
$1145$ at $123$ operations and the previous $116$-operation degree $1353$.
The central improvement is to shift the index quotient by $q$ and replace a
large radical-separation bound by the two-operation expression $I(T+I)$.
We also combine the odd-scale theorem with the existing integer lifts to
obtain a quartic in $51$ unrestricted integers, a sextic in $42$, and an octic
in $36$, reducing each count by three. A separate construction gives an
explicit equation in $17$ unrestricted integer unknowns, compared with the
$26$ in the reviewed direct compiler family. It uses $20\,607$ operations
and has exact degree $424296\cdot5^{33}+53824$; this is a variable-count
improvement with a very large degree cost.
The coefficients for the twelve-variable and low-degree forms are the complete
materialized $943$-tile tuple, reproduced byte for byte. The seventeen-variable
form retains the entire original package-18 T4 tuple and its finite index recipe.
For each stated form, ordinary mathematics formalizable in ZFC proves that
absence of solutions is equivalent to $\Con(\mathrm{ZFC})$. The inherited
compiler and logical calibration are identified explicitly; finite audits are
not represented as new proof-assistant certificates. No absolute minimality
or proof of consistency is claimed.
\end{abstract}
\clearpage
\tableofcontents
\newpage

\section{The result and the comparison standard}
\label{sec:results}

The source repository was inspected at the immutable revision
\begin{center}\small\pin\end{center}
Its current catalogue contains eight wave-3 packages, numbered 16--23; earlier
waves survive in history and in copied references. We compare against that
catalogue, rather than an older thirteen-variable starting point
\cite{catalogue}. A smaller equation can mean fewer unknowns, lower degree,
fewer operations, smaller coefficient height, or a simpler proof. These
criteria disagree. Table~\ref{tab:frontier} records the improvements relevant
to the present construction without collapsing them into one ordering.

\begin{table}[ht]
\centering\small
\begin{tabular}{@{}lrrrr@{}}
\toprule
Construction and domain & Unknowns & Equations & Exact degree & Operations\\
\midrule
Package 20 product, $\Pos$ &12&1&1145&123\\
Package 21 product, $\Pos$ &12&1&1353&116\\
Package 21 system, $\Pos$ &12&3&1344&112\\
Package 22 system, $\Pos$ &12&3&1520&112\\
\textbf{New product}, $\Pos$ &\textbf{12}&\textbf{1}&\textbf{1097}&\textbf{116}\\
\textbf{New system}, $\Pos$ &\textbf{12}&\textbf{3}&\textbf{1088}&\textbf{110}\\
New sum of squares, $\Pos$ &12&1&2176&118\\
\midrule
Package 20 quartic, $\Z$ &54&1&4&--\\
\textbf{New quartic}, $\Z$ &\textbf{51}&\textbf{1}&\textbf{4}&--\\
Package 20 sextic, $\Z$ &45&1&6&--\\
\textbf{New sextic}, $\Z$ &\textbf{42}&\textbf{1}&\textbf{6}&--\\
Package 20 octic, $\Z$ &39&1&8&--\\
\textbf{New octic}, $\Z$ &\textbf{36}&\textbf{1}&\textbf{8}&--\\
\bottomrule
\end{tabular}
\caption{Degrees of systems mean maximum residual degree. Operation counts
charge every binary $+,-,\times$, including multiplication by numerals;
reuse and equality comparisons are free. A dash means no optimized arithmetic
operation count is asserted.}\label{tab:frontier}
\end{table}

The new product lowers the best reviewed twelve-variable degree by $48$ and
uses seven fewer operations than that degree-minimizing predecessor. Against
the old $116$-operation product, it lowers degree by $256$. The $110$-operation
system uses a shifted transport coordinate and compares the two sides of the
last norm directly; these implementation choices are disclosed in
Section~\ref{sec:systems}. Thus the count is not obtained by silently changing
the operation model.

Several existing forms remain preferable under other criteria. The explicit
nine-natural-variable Jones specialization in package 18 still uses fewer
unknowns than the moderate-degree construction. Its degree is
$161844\cdot5^{33}+484$ and its delivered circuit uses $20\,531$ operations.
The older $67$-operation product still has the smallest reported single-equation
operation count. The $59$-operation system in sixteen positive variables is
also unaffected. No claim below beats those counts \cite{catalogue,jones9}.

For unrestricted integers, Section~\ref{sec:integer17} gives a fully specified
seventeen-variable equation. Its exact degree is
$424296\cdot5^{33}+53824$ and its circuit has $20\,607$ binary operations.
It improves the concrete $26$-integer count in the inspected direct family.
It is not a global variable-count record: Sun's 2026 representation theorem
gives a general ten-integer bound \cite{sun2026}. This article supplies no
concrete ten-integer specialization and makes no claim of absolute minimality.

\paragraph{Meaning of the logical assertion.}
Fix the syntactic proof predicate used by the K215 endpoint. Writing
$\mathrm{Ref}_{\mathrm{ZFC}}$ for existence of a finite ZFC derivation of
falsity, the theorem asserted for each construction is
\[
 \mathrm{ZFC}\vdash
 \bigl(\exists\mathbf x\in\mathcal D:\ P(\mathbf x)=0\bigr)
 \longleftrightarrow\mathrm{Ref}_{\mathrm{ZFC}}.
\]
Negation gives the requested equivalence with $\Con(\mathrm{ZFC})$.
This does not prove that either side of the negated equivalence holds.
No additional claim that ZFC cannot prove inconsistency is made under
mere consistency. The domain $\mathcal D$ is part of each theorem.

\section{Fixed coefficients and inherited mathematical dependencies}
\label{sec:coefficients}

The twelve-variable and low-degree constructions use the complete
\path{wang943-compact-5pow10} tuple from package 19 \cite{wang943}. Put
\[
 d_\ast=5^{10}=9\,765\,625,\qquad B=2^{d_\ast},\qquad
 \beta=B-1,\quad K_0=D_C+BD_R.
\]
The remaining coefficients are $M_C$ and the \emph{shifted} source mask
$M_F=M_F^{\mathrm{native}}+B-1$. In particular the source port $M_F$
must not be replaced by the native mask. Every coefficient is fixed before
any unknown is quantified. No coefficient depends on whether ZFC is consistent.

The archive includes the literal K215 transition table, the finite coefficient
generators, the materialized coefficient payloads, and their hashes. The
coefficient sizes are:
\begin{center}
\begin{tabular}{@{}lrr@{}}\toprule
Port & Bits & Uncompressed bytes\\\midrule
$\beta$&9\,765\,625&1\,220\,704\\
$K_0$&13\,201\,251&1\,650\,157\\
$M_C$&9\,765\,625&1\,220\,704\\
$M_F$&9\,765\,626&1\,220\,704\\\bottomrule
\end{tabular}
\end{center}
The four deterministic gzip payloads total $20\,260$ bytes. This is a storage
fact, not a small coefficient-height claim. All four reproduced raw SHA-256
digests agree with the pinned manifest.

For precision, the bounded recipe uses $\rho=2^{625}$ and a cell length
$L_c=15625$. Its $943$ ordered directed-signal Wang tiles are generated from
all $430$ transitions of K215. Selector 0 is the true blank-start tile;
selector 1 is a globally impossible tile. Binary occupied edge labels produce
forty copy positions and four anchors. If $\mathcal E$ is the resulting
$989$-position support and $c_e$ its explicitly generated clause coefficients,
the compact geometry is
\[
 H_c=5497,\quad T_1=7631,\quad T_2=8733,
 \quad\text{anchors }1095,1096,2195,3296.
\]
For $\kappa=1024$, the clause mask is
\[
 \mu=\kappa-2+\sum_{j=1}^{44}(\kappa-3)\kappa^j
       +2^{450}(2^{78}-1).
\]
The integers are fixed by
\begin{align*}
 D_C={}&\rho^{30}+\rho^{5517}+\rho^{1100}+\rho^{2200}
 +\sum_{e\in\mathcal E}c_e(\rho^{7631-e}+\rho^{8733-e}),\\
 D_R={}&\rho^{5497},\\
 M_C={}&B-1-\sum_{e\in\mathcal E\setminus\{1\}}\rho^e-2\rho^{1063},\\
 M_F^{\mathrm{native}}={}&
 \mu(\rho^{7631}+\rho^{8733})
 +\sum_{j=0}^{9}\bigl(\rho^{1013+j}+\rho^{6520+j}\bigr)
 +\rho^{2195}+\rho^{3296}+4.
\end{align*}
The complete finite rule for $\mathcal E,c_e$ and the tile ordering is in
\texttt{coefficients/build\_wang943.py}; the compact geometric substitution
is in \texttt{build\_wang943\_compact.py}. Thus this is a particular
integer polynomial, not a universal polynomial with an unspecified program
index. The accompanying source proof establishes
\[
 0<M_C,M_F^{\mathrm{native}}<B-1,\quad
 M_C\equiv2\pmod4,\quad M_F^{\mathrm{native}}\equiv4\pmod8,
\]
and $\pc(M_C)+\pc(M_F^{\mathrm{native}})=d_\ast$.

\subsection{What is imported and what is proved here}

The logical chain has three components.
First, the literal blank run of K215 halts if and only if ZFC has a finite
refutation. The pinned Lean declaration
\texttt{KEndpoint215.blank\_halts\_iff\_zfc\_refutation} states that
calibration; its source is included for identification \cite{k215}.
Second, the native-window/Pell compiler translates that blank halting
predicate into the arithmetic interface recalled below. Its ordinary
mathematical proof, including the Wang adapter, is an inherited dependency.
Third, the transformations proved in this article replace that interface
by the displayed polynomials while preserving existence of solutions.

The first two components are not reproved from axioms here. Relevant source
proofs are copied into \texttt{references/} and \texttt{coefficients/}.
The new equivalence proof is therefore modular: it states and uses the
precise existing interface, rather than treating a finite circuit test as
proof of the compiler. No Lean or Rocq build of this new development was run.
The assertion of provability in ZFC comes from formalizing the ordinary
finite arithmetic and Pell arguments in ZFC, not from an assumed implication
from Lean's metatheoretic soundness to ZFC.

\section{The twelve-variable polynomial}
\label{sec:definition}

All twelve supplied unknowns are strictly positive integers:
\[
 (J,F,\alpha,z,h,i,T,s,w,\eta,\zeta,r)\in\Pos^{12}.
\]
Every other symbol in this section is a polynomial abbreviation, not an
additional unknown. Define
\begin{align*}
 q&=\beta J+1,& C&=q-F-\alpha,&X&=wq,&Y&=sq^3,& E&=XY,\\
 k&=\eta+\zeta,&c&=kY+\eta,&a&=E+Y,&H&=4a+3,&\Delta&=a^2+H,\\
 \ell&=EkY,&\mathcal A&=1+\ell(\ell+k),&G&=1+\Delta c^2.
\end{align*}
The packed index and the two outer factors are
\begin{align}
 R&=\bigl((q-1)(q-F)+\alpha\bigr)(q^2-1)+(M_C+qM_F)J,\label{eq:R}\\
 N_k&=k-(q+h)E-R,\label{eq:Nk}\\
 N_t&=(K_0+w)C+q-F-z(q-1).\label{eq:Nt}
\end{align}
The term $q+h$ replaces the old positive index quotient. This is an actual
coordinate restriction whose completeness will be proved, not an identity
on arbitrary old witnesses.

Define the strong and auxiliary expressions
\begin{align*}
 I&=ic^2,&D&=1+\Delta I^2,&Q&=(\Delta I)^2,&d&=Q-1,\\
 v&=cT,&b&=c+RD,&e&=Q(Dv^2+b^2)-1,&B_a&=2Qvb.
\end{align*}
The new, much shorter separation bound is
\begin{equation}
 \boxed{\quad W=I(T+I),\qquad L=1+HW,\qquad S=X+ac+Hr.\quad}
 \label{eq:weight}
\end{equation}

Temporarily use formal square roots $u^2=G$, $t^2=\mathcal A$, and
$f^2=D$. Define
\begin{equation}
 \mathcal N=
 \prod_{\sigma,\tau,\upsilon\in\{-1,1\}}
 \left[d(S-\sigma Lu-\tau HLt)^2-H^2e
              +\upsilon H^2B_af\right].\label{eq:norm}
\end{equation}
The product is invariant under the three formal sign changes, so every
square root cancels and $\mathcal N$ is a polynomial with integer
coefficients in the twelve supplied variables. This statement holds even
when the specialized radicands have dependent square classes or are squares.
The radicals in \eqref{eq:norm} are notation for a finite polynomial
identity; there are no hidden existential root variables.

The requested single equation is
\begin{equation}
 \boxed{\qquad P_{12}=N_kN_t(1+\mathcal N)-1=0.\qquad}
 \label{eq:main}
\end{equation}
Section~\ref{sec:circuit} gives a radical-free evaluation of
$\mathcal N$. The file \texttt{norm12/shifted\_norm12\_product.json}
lists every one of its $116$ binary operations.

\begin{theorem}\label{thm:main}
With the fixed coefficient tuple of Section~\ref{sec:coefficients}, ZFC
proves that \eqref{eq:main} has a positive integer solution if and only
if K215 halts on its blank tape, and hence if and only if ZFC has a finite
refutation. Equivalently,
\[
 \mathrm{ZFC}\vdash
 \left(\nexists\mathbf x\in\Pos^{12}:P_{12}(\mathbf x)=0\right)
 \longleftrightarrow\Con(\mathrm{ZFC}).
\]
The polynomial has exact degree $1097$ and a $116$-operation circuit.
\end{theorem}

\section{Bounds established before any root is recovered}
\label{sec:bootstrap}

The parity argument in Section~\ref{sec:parity} shows that a zero of
$P_{12}$ satisfies $\mathcal N=0$ and $N_k=N_t=\epsilon$ for one
$\epsilon\in\{-1,1\}$. We first derive all inequalities used by the
new root bound from these outer unit equations alone.

If $C\leq-1$, then
\[
 q-F-z(q-1)=1-F-(z-1)(q-1)\leq0
\]
and $N_t\leq-(K_0+w)\leq-2$, a contradiction. Thus $C\geq0$;
consequently $q-F=C+\alpha\geq1$ and $1\leq F\leq q-1$.
Let $G_0=(q-1)(q-F)+\alpha=q\alpha+(q-1)C$. Then
\[
 q\leq G_0\leq q(q-1),\qquad
 0<(M_C+qM_F)J<(1+2q)(q-1).
\]
It follows directly from \eqref{eq:R} that
\begin{equation}
 q(q^2-1)<R<q^4-q^3+q^2-1<q^4-2.\label{eq:Rbounds}
\end{equation}
Also $q=\beta J+1\geq B>16$, $R>3q$, $Y\geq q^3$, and
$E\geq q^4>R+2$. The shifted quotient $q+h$ is positive independently
of the value of $R$.

The useful effect of the shift is the inequality $c>R(a+5)$.
Indeed
\begin{align*}
 (q+1)EY-q^4(E+Y)
 &=qE(Y-q^3)+Y(E-q^4)\geq0.
\end{align*}
Since $R<q^4$ and $a=E+Y$, this gives $(q+1)EY>Ra$. From
$N_k=\epsilon$ and $h\geq1$ we obtain
\[
 c=(q+h)EY+(R+\epsilon)Y+\eta
    >Ra+(R-1)Y>R(a+5),
\]
where the last step uses $R\geq2$ and $Y\geq4096$.
As $a\geq5$, we have $\Delta=a^2+4a+3<a(a+5)$. Therefore
\begin{equation}
 \boxed{\quad ac>R\Delta,\qquad c>a,\qquad c>R.\quad}
 \label{eq:keygap}
\end{equation}
These facts do not assume that $u,t,f$, or an auxiliary root is rational.
This order is essential: using a Pell classification to justify
\eqref{eq:keygap} here would be circular.

\section{The compact separation bound and recovery of all roots}
\label{sec:separation}

If $\mathcal N=0$, choose one vanishing factor in \eqref{eq:norm} and
absorb its three signs into coherent choices of $u,t,f$. Set
$y=(S-Lu-HLt)/H$. Then
\begin{align}
 u^2&=G,&t^2&=\mathcal A,&f^2&=D,\notag\\
 d y^2&=e-B_af=Q(vf-b)^2-1,&
 S&=Lu+HLt+Hy.\label{eq:coherent}
\end{align}
Work in a finite normal extension of $\mathbb Q$ containing these roots.
Every conjugate $y'$ has an associated $f'=\pm\sqrt D$. Hence
\begin{equation}
 |y'|<\sqrt{Q/d}\,(v\sqrt D+b).\label{eq:rawbound}
\end{equation}
This includes complex conjugates. To see it explicitly, put
$M=v\sqrt D+b>1$. The real number $(vf'-b)^2$ is between $0$ and
$M^2$, so the absolute value of $Q(vf'-b)^2-1$ is at most
$\max(1,QM^2-1)<QM^2$.

\begin{lemma}[The two exact square comparisons]\label{lem:bound}
Under the outer unit equations and positive supplied coordinates,
\[
 \sqrt{Q/d}\,c\sqrt D<I\sqrt G,
 \qquad \sqrt{Q/d}\,b<I^2\sqrt G.
\]
Consequently every conjugate in \eqref{eq:coherent} satisfies
$|y'|<I(T+I)\sqrt G=W\sqrt G$.
\end{lemma}
\begin{proof}
For the first inequality, an exact expansion gives
\[
 dI^2G-Qc^2D
  =I^2\bigl[\Delta^2(I^2-c^2)-G\bigr]>0.
\]
Indeed $I\geq c^2$, $c\geq2$, and $\Delta\geq2$, so the bracket
is at least $\Delta(\Delta-1)c^2-1>0$.

For the second inequality, \eqref{eq:keygap} and integrality give
$ac-R\Delta\geq1$. Also $c+R<2c\leq c^4\leq I^2$. Thus
\[
 b=c+R+R\Delta I^2<acI^2.
\]
The remaining comparison is
\[
 dG-Qa^2c^2=Q(1+Hc^2)-G
           =(Q-1)+(QH-\Delta)c^2>0.
\]
Therefore $\sqrt{Q/d}\,ac<\sqrt G$, proving the second inequality.
Multiplying the first by $T>0$ and adding the second in
\eqref{eq:rawbound} gives the stated bound.
\end{proof}

Because $L=1+HW$, the lemma implies
\begin{equation}
 H|y'|<L\sqrt G.\label{eq:domu}
\end{equation}
Moreover $\sqrt{\mathcal A}>\ell=kEY$, while
$\sqrt G<(a+2)c<2(a+2)kY$. Since $HE>4(a+2)$,
\begin{equation}
 H\sqrt{\mathcal A}>2\sqrt G.\label{eq:domt}
\end{equation}

If $t$ were irrational, an automorphism of the normal extension would
send $t$ to $-t$. Subtract its image of the last equality in
\eqref{eq:coherent}. The changed $HLt$ term has magnitude
$2HL\sqrt{\mathcal A}$, whereas the combined magnitude of the other
changed terms is less than $4L\sqrt G$, by \eqref{eq:domu}.
This contradicts \eqref{eq:domt}. Thus $t$ is rational, and therefore
integral because its square is an integer.

Now an automorphism negating an irrational $u$ would fix $t$. The same
subtraction would give $2L\sqrt G<2L\sqrt G$, again impossible.
Thus $u$ is an integer. Equation~\eqref{eq:coherent} makes $y$ rational.
Since $B_a>0$, the equality $f=(e-dy^2)/B_a$ makes $f$ rational and
then integral. The last integrality step for $y$ needs a special Pell
lemma; rationality alone is insufficient.

\begin{lemma}[Rational auxiliary Pell roots]\label{lem:saturation}
If $m\geq2$ and $V$ are integers, $y\in\mathbb Q$, and
$m^2V^2-(m^2-1)y^2=1$, then $y\in\mathbb Z$.
\end{lemma}
\begin{proof}
Write $y=A_1/B_1$ in lowest terms with $B_1>0$. The equation gives
$B_1^2\mid m^2-1$. Put $D_1=(m^2-1)/B_1^2$. This is a positive
nonsquare integer. The two positive Pell solutions are
$(m,B_1)$ and $(m|V|,|A_1|)$; neither ordinate is zero.
If $\chi_j+\psi_j\sqrt{D_1}$ denotes the $j$th positive power of
the fundamental positive Pell unit with integer coefficients in
$\Z[\sqrt{D_1}]$, write these solutions as
$(\chi_j,\psi_j)$ and $(\chi_l,\psi_l)$.

The elementary divisibility fact $\chi_j\mid\chi_l$ forces $l$ to
be an odd multiple of $j$. To verify it, the Pell addition formulas give
$\chi_{n+2j}\equiv-\chi_n\pmod{\chi_j}$ and reflection reduces
the index modulo $2j$ into $0\leq n\leq j$. For $n<j$, strict
growth gives $0<\chi_n<\chi_j$. The only divisible case is $n=j$.
Pell multiplication then gives $\psi_j\mid\psi_l$, so
$B_1\mid A_1$. Coprimality forces $B_1=1$.
\end{proof}

Apply the lemma to $m=\Delta I$ and $V=vf-b$. It proves that $y$
is integral. Finally $L\equiv1\pmod H$ in \eqref{eq:coherent} gives
$u\equiv X+ac\pmod H$, so
\[
 \gamma=(u-X-ac)/H\in\mathbb Z.
\]
The restored coordinates satisfy the first, main, strong, and auxiliary
Pell equations of the inherited signed normalized interface:
\begin{align*}
 t^2-\ell(\ell+k)&=1,\\
 (X+ac+H\gamma)^2-\Delta c^2&=1,\\
 f^2-\Delta i^2c^4&=1,\\
 Q\bigl((c(Tf-1)-Rf^2)^2-y^2\bigr)+y^2&=1.
\end{align*}
The auxiliary identity uses $f^2=D$. The inherited signed-rank theorem
resolves the common outer unit sign to $+1$ and decodes an actual blank
halting computation. Signed recovered roots and $\gamma$ are permitted
by that interface; no positivity of those recovered coordinates is assumed
\cite{norm12,blank67}.

\section{Completeness of the quotient shift}
\label{sec:complete}

It remains to prove that imposing $h_{\mathrm{old}}=q+h$ loses no
accepting witnesses. Take a canonical positive solution of the old blank
compiler. The existing first-Pell classification gives
\[
 k=2\psi_{P_0}(n),\qquad P_0=2EY+1,
 \qquad 2n=R+1.
\]
Here $\psi_P(0)=0$, $\psi_P(1)=1$, and
$\psi_P(n+2)=2P\psi_P(n+1)-\psi_P(n)$. The equality $2n=R+1$
is a conclusion of the complete old signed-rank argument, not a premise
used in Section~\ref{sec:bootstrap} \cite[Section 4]{blank67}.
As $R>3q$, we have $n\geq3$, whence
\[
 k\geq2\psi_{P_0}(3)=8P_0^2-2>32E^2Y^2.
\]
The old index equation gives
\[
 h_{\mathrm{old}}=(k-R-1)/E>R>q,
\]
using $R+1<E$ and $R<E$. Therefore
$h=h_{\mathrm{old}}-q$ is a positive integer.

For the canonical positive roots $u,t,f,y$ and $\gamma$, define
\[
 r=\gamma+Wu+Lt+y>0.
\]
The identity $u=X+ac+H\gamma$ gives
$S=Lu+HLt+Hy$. One factor of \eqref{eq:norm} vanishes, so
$\mathcal N=0$ and $P_{12}=0$. This proves completeness. Together with
the soundness argument above and the calibration of K215 it proves the
semantic part of Theorem~\ref{thm:main}.

\section{A radical-free circuit, the parity guard, and exact degree}
\label{sec:circuit}

For an explicit finite formula with no displayed radicals, use pairs
$(x_0,x_1)$ representing $x_0+x_1u$ in
$\Z[\mathbf x][u]/(u^2-G)$. Pair multiplication is
\[
 (x_0,x_1)(y_0,y_1)=(x_0y_0+Gx_1y_1,\ x_0y_1+x_1y_0).
\]
All pair operations below are this fixed finite arithmetic rule. Put
\begin{align*}
 A_0&=H^2L^2\mathcal A,&V&=(S^2+L^2G,\ 2SL),\\
 \mathcal C=C_0+C_1u&=d(V+A_0)-H^2e,&K&=4d^2A_0V,\\
 Z_0&=H^4B_a^2D,&N_0+N_1u&=(\mathcal C^2-K-Z_0)^2-4KZ_0.
\end{align*}
Scalars such as $A_0,H^2e,Z_0$ are inserted as pairs with second
coordinate zero. Then
\begin{equation}
 \mathcal N=N_0^2-GN_1^2.\label{eq:pairnorm}
\end{equation}
Expanding the two conjugates in $t$ and the two in $f$ verifies this
identity with \eqref{eq:norm}. The identity
\[
 (\mathcal C^2+K-Z_0)^2-4\mathcal C^2K
 =(\mathcal C^2-K-Z_0)^2-4KZ_0
\]
avoids a full pair multiplication because $Z_0$ is scalar. The circuit
also uses $Z_0=4H^4(QDv^2)(Qb^2)$, reusing products already needed
for $e$. These optimizations are inherited from package 21; the new
changes are the quotient shift and smaller bound.

\subsection{Why the product really enforces all equations}
\label{sec:parity}

The coefficient of $u$ in $V$ is even; hence $C_1$ is even, the odd
coefficient of $\mathcal C^2$ is divisible by $4$, and that of $K$ by $8$.
The odd coefficient of $(\mathcal C^2-K-Z_0)^2$ is therefore divisible by $8$,
as is that of $4KZ_0$. Thus $8\mid N_1$ as an integer polynomial and
\[
 \mathcal N\equiv N_0^2\pmod{64}.
\]
In particular $\mathcal N\ne-2$. If
$N_kN_t(1+\mathcal N)=1$, all three factors are integer units.
The last cannot be $-1$, so it is $1$, $\mathcal N=0$, and $N_k=N_t$
are common units. This proves the assertion used at the start of
Section~\ref{sec:bootstrap}, without any sign or root assumption.

\subsection{Uniform exact degree}

Treat the four fixed ports as degree zero and every supplied unknown as
degree one. The important degrees are
\begin{center}
\begin{tabular}{@{}lr@{\qquad}lr@{\qquad}lr@{}}\toprule
Expression&Degree&Expression&Degree&Expression&Degree\\\midrule
$c$&5&$a,H,E$&6&$\Delta$&12\\
$\mathcal A,G$&22&$I$&11&$D$&34\\
$Q,d$&46&$b$&38&$e$&122\\
$W$&22&$L$&28&$S$&11\\
$A_0$&90&$C_0$&136&$C_1$&85\\
$N_0$&544&$N_1$&$\leq493$&$\mathcal N$&1088\\\bottomrule
\end{tabular}
\end{center}
The highest term of $C_0$ comes uniquely from $dA_0$: the competing
$H^2e$ has degree $134$ and $dV_0$ degree $124$. Therefore the leading
term of $\mathcal N$ comes from $C_0^8$ and cannot cancel against
$GN_1^2$, whose degree is at most $1008$.

For a uniform certificate put $q_0=\beta J$ and $k_0=\eta+\zeta$.
The leading homogeneous form is
\[
 \mathcal N_{\mathrm{top}}=
 2^{64}i^{48}k_0^{112}w^{80}s^{192}q_0^{656}.
\]
The leading factors of $N_k$ and $N_t$ are respectively
\[
 -(q_0+h)wsq_0^4,
 \qquad w(q_0-F-\alpha)-zq_0.
\]
Thus $\deg P_{12}=1088+7+2=1097$. For example the coefficient of
\[
 J^{661}hz\,i^{48}\eta^{112}w^{81}s^{193}
\]
is $2^{64}\beta^{661}>0$. This proves exact degree for the authentic
coefficient tuple, not just for a small numerical specialization.

\section{The 110-operation system and domain conventions}
\label{sec:systems}

The unshifted three-equation presentation is
$N_k=1$, $N_t=1$, and $\mathcal N=0$. Its most economical delivered
comparison form uses $111$ operations, since \eqref{eq:pairnorm} can be
checked as $N_0^2=GN_1^2$ without calculating their difference.

For one further saving substitute $z_{\mathrm{old}}=z+1$ in the
transport equation. The result is
\begin{equation}
 \boxed{\quad N_k=1,\qquad (K_0+w)C=F+z(q-1),\qquad
                  N_0^2=GN_1^2.\quad}\label{eq:system}
\end{equation}
The norm is independent of $z$. The exact identity underlying the shift is
\[
 (K_0+w)C+q-F-(z+1)(q-1)-1
 =(K_0+w)C-F-z(q-1).
\]
Every positive shifted solution gives a positive unshifted solution.
Conversely a canonical old witness has $C>0$, $X=wq=2^R$, and $R>3q$,
so $w>q>F$. The old transport equation then forces
$(z_{\mathrm{old}}-1)(q-1)=(K_0+w)C-F>0$; hence
$z_{\mathrm{old}}\geq2$. The shift is complete.

The exact operation ledgers are
\begin{center}
\begin{tabular}{@{}lrrr@{}}\toprule
Form & Multiplications & Additions/subtractions & Total\\\midrule
Product $P_{12}$ &73&43&116\\
System \eqref{eq:system} &71&39&110\\
Sum of the three squared residuals &74&44&118\\\bottomrule
\end{tabular}
\end{center}
All gates and all supplied ports are live. The sum of squares has degree
$2176$, since $\mathcal N^2$ is its unique highest-degree summand.
The product uses the old transport coordinate; the system and its sum of
squares use the shifted one. They have equivalent solvability predicates,
not identical named-coordinate zero sets.

For a formulation over $\N=\{0,1,2,\ldots\}$, replace each of the
twelve positive inputs by one plus a new natural input. This preserves
variable count and degree and adds at most twelve input additions to the
literal circuit. Reading the displayed positive-domain polynomial directly
over all integers is not justified. No $26$-integer extension of the new
compact bound is asserted.

% Input by report.tex. All mathematical abbreviations in this section are local.
\section{A concrete equation in seventeen unrestricted integers}
\label{sec:integer17}

This section gives a fixed polynomial $P_{17}$ in seventeen unrestricted
integer unknowns whose absence of zeros is equivalent, provably in ZFC, to
$\Con(\mathrm{ZFC})$. It uses the complete original T4 coefficient tuple of
package 18, \texttt{18-jones9}, at revision \pin\ \cite{jones9}.
All source coefficients and abbreviations below are local to this section.
In particular, this construction has its own coefficient binding, distinct
from the $943$-tile binding used in the preceding sections.

The result improves the delivered concrete $26$-integer count in the
inspected direct compiler family, at the cost of an enormous degree.
It establishes neither global minimality nor a universal-variable record:
Sun's ten-integer representation theorem already gives a smaller abstract
general bound \cite{sun2026}. The contribution here is the complete fixed
specialization, its explicit circuit, and the additional proofs needed to
use signed coordinates. The source Jones ratio theorem, quartic/compiler
theorem, and K215 calibration are inherited ordinary mathematical theorems.
The new arguments below are proved in full; no new proof-assistant
certificate is claimed.

\subsection{The fixed coefficients and exact source interface}
\label{i17:source}

Start with eight base coordinates
\[
 (\varepsilon,g_0,h_0,s_0,w_0,\phi,i_0,j_0)\in\Z^8
\]
and abbreviate
\[
 g=g_0+1,\quad h=h_0+1,\quad s=s_0+1,\quad
 w=w_0+1,\quad i=i_0+1,\quad j=j_0+1.
\]
The old ninth coordinate of package 18 and its finalizer are replaced
below. The full source expressions through the ratio conditions are
reproduced here so that the signed-domain argument has an explicit
interface.

\paragraph{Finite coefficient recipe.}
Put $\nu=32$ and $L=5^{33}$. Shift the $32$ positive coordinates of the
pinned quartic by $1$, and include its unused dummy input as coordinate
zero. The dummy is fixed to $1$ and is not an unknown of $P_{17}$.
Let $p_\alpha$ be the resulting coefficient indexed by
$\alpha\in\N^{33}$, with
$|\alpha|=\sum_{r=0}^{32}\alpha_r\le4$, and put
\[
 c_\alpha=(4-|\alpha|)!\left(\prod_{r=0}^{32}\alpha_r!\right)p_\alpha.
\]
The four compiler ports are bound to the complete original T4 tuple in
\path{signed17/upstream/COEFFICIENT_BINDINGS.json}.
The coefficient records and the positive-coefficient bounds $\rho_j$ for
the shifted source residuals are specified by
\path{signed17/upstream/nine/generated/constants.json} and
\path{signed17/upstream/nine/generated/index_bound_s.json}.
Define the following fixed integers and fixed polynomial:
\begin{align*}
 s_\star&=2+24\sum_j\rho_j^2,
 &z_\star&=2^{s_\star},\\
 e_\star(T)&=
 \sum_{\substack{\alpha\in\N^{33}\\|\alpha|\le4}}
 (z_\star+c_\alpha)T^{\,L-\sum_{r=0}^{32}\alpha_r5^r},
 &y_\star&=e_\star(2z_\star).
\end{align*}
The sum includes the unlisted zero coefficients. All exponents are fixed
nonnegative integers. The source bound proves
\[
 z_\star>|c_\alpha|+1\quad\text{for every }\alpha,
 \qquad y_\star\ge2z_\star.
\]
Thus the coefficient recipe is fixed independently of the quantified
unknowns and of the consistency of ZFC.

\paragraph{Source polynomial abbreviations.}
Define
\begin{align*}
 b&=y_\star+1+\varepsilon,
 &\mathcal B&=b^5,
 &q&=\mathcal B^L,\\
 \ell&=\sum_{r=1}^{\nu}\mathcal B^{5^r},
 &\lambda&=\sum_{r=0}^{4L-1}\mathcal B^r,
 &e&=e_\star(\mathcal B),\\
 \theta&=\mathcal B-2z_\star,
 &N&=q^{16},
 &G&=32b\mathcal B^{5^\nu}.
\end{align*}
To display the packed source expression without a long line, write
\begin{align*}
 R_1={}&g+eq^3+\ell q^5\\
 &+\bigl[2(e-z_\star\lambda)(1+\mathcal B+g)^4
       +\lambda\mathcal B(1+q^4)\bigr]q^7,\\
 R_2={}&q^3-b\ell+\ell+\theta\lambda q^3+(\mathcal B-2)q^8,\\
 R={}&R_1(N^2-N)+R_2(N^2-1).
\end{align*}
The remaining abbreviations are
\begin{align*}
 U&=N^2w,&Y&=N^2s,&M&=RY,\\
 P&=2M^2U,&K&=R+1+h(P-1),&A&=M(U+1),\\
 B_1&=2R+1,&W_1&=bw,&C&=B_1+W_1+\phi,\\
 J&=4A-5,&Q&=3W_1C-2(W_1^2-1),
\end{align*}
and
\begin{align*}
 D&=(A^2-1)C^2+1,&E&=iJDC^2,\\
 F&=(A^2-1)E^2+1,&G_1&=A+F(CD+1-A),\\
 H&=B_1+jC,&I&=(G_1^2-1)H^2+1,\\
 a_0&=(P^2-1)K^2+1,&a_1&=DFI,\\
 D_0&=JF,&T_0&=QF+J(H-C).
\end{align*}
Finally, set
\[
 m=K^2G-4(C-KY)^2G-8gK^2.
\]
Let $\operatorname{Sq}(\Z)=\{r^2:r\in\Z\}$.
The imported Jones ratio equivalence states that there exist eight
\emph{nonnegative} base coordinates satisfying
\begin{equation}
 a_0,a_1\in\operatorname{Sq}(\Z),\qquad
 D_0\mid T_0,\qquad m>0
 \label{i17:JR}
\end{equation}
if and only if the pinned $32$-variable quartic has a positive solution.
The imported compiler and K215 calibration identify that assertion with
existence of a finite ZFC refutation. These are the precise inherited
theorem boundaries used below. The signed-coordinate argument depends
on the displayed source structure, not merely on existence of an
unspecified nine-variable representation.

\subsection{Two nonnegativity guards are redundant}

Only six base coordinates need explicit nonnegativity guards:
\[
 \mathbf x=(x_1,\ldots,x_6)
          =(\varepsilon,g_0,s_0,w_0,\phi,j_0).
\]

\begin{lemma}[The source denominator never vanishes]
\label{i17:denominator}
For every integer assignment to the eight base coordinates, $D_0\ne0$.
\end{lemma}
\begin{proof}
The integer $J=4A-5$ cannot vanish. If $F=0$, then
$(A^2-1)E^2=-1$, which forces $A=0$ and $E=\pm1$.
But $A=0$ gives $J=-5$, so $E=iJDC^2$ is divisible by $5$, a contradiction.
\end{proof}

\begin{lemma}[A signed coordinate]
\label{i17:signed-i}
If $i\ne0$, replacing $i$ by $|i|$ leaves
$a_0,a_1,D_0,T_0,m$ unchanged.
\end{lemma}
\begin{proof}
The coordinate $i$ enters through $E$, and these five expressions use
$E$ only through $E^2$ in $F$. The replacement therefore preserves all
five values. The new base coordinate $i_0=|i|-1$ is nonnegative.
\end{proof}

\begin{lemma}[The guard on $h_0$ follows from the source conditions]
\label{i17:h-guard}
Suppose that the six coordinates $\mathbf x$ are nonnegative,
$a_0\in\operatorname{Sq}(\Z)$, and $m>0$. Then $h\ge1$, so $h_0\ge0$.
\end{lemma}
\begin{proof}
The coefficient recipe gives $b\ge2$, $\mathcal B=b^5\ge32$,
$e\ge0$, and $\mathcal B>2z_\star$. Put $t=5^\nu$.
In particular, $G>0$ and $g\ge1$. The positive margin first implies
\begin{equation}
 K\ne0,\qquad g<G/8=4b\mathcal B^t,\qquad
 4(C-KY)^2<K^2.
 \label{i17:margin-bounds}
\end{equation}
Indeed, $K=0$ would make $m\le0$, while
$m\le(G-8g)K^2$ gives the bound on $g$.
Deleting the strictly positive term $8gK^2$ from the margin inequality
and dividing by $G$ gives the last inequality.

Since $b\le\mathcal B$ and $\mathcal B>5$, we have
\begin{equation}
 \begin{aligned}
 1+\mathcal B+g
 &<1+\mathcal B+4\mathcal B^{t+1}
 \le5\mathcal B^{t+1}\\
 &<\mathcal B^{t+2}\le\mathcal B^{5t}=q.
 \end{aligned}
 \label{i17:g-bound}
\end{equation}
The distinct exponents in $\ell$ lie between $0$ and $t$, so
\[
 \ell\le\sum_{r=0}^t\mathcal B^r<\mathcal B^{t+1},
 \qquad b\ell<\mathcal B^{t+2}\le q<q^3.
\]
Writing $V_0=1+\mathcal B+g$, the possibly negative part of $R_1$ is
dominated by its positive companion:
\[
 \lambda\mathcal Bq^4-2z_\star\lambda V_0^4>0.
\]
This follows from \eqref{i17:g-bound} and $\mathcal B>2z_\star$.
After separating $2(e-z_\star\lambda)V_0^4$ into its two terms,
all remaining terms of $R_1$ are nonnegative, and $g>0$.
Also $R_2>0$, since $q^3>b\ell$, $\theta>0$, and $\mathcal B>2$.
Consequently
\[
 R\ge N^2-N\ge2.
\]
This proves positivity of $R$ from the margin-derived bound, without
assuming any sign for $h$.

Now $Y,U\ge1$, $C=2R+1+bw+\phi>0$, and
$P=2R^2Y^2U\ge2R^2>R+2$.
If $K<0$, then
$|C-KY|=C+|K|Y>|K|$, contrary to \eqref{i17:margin-bounds}.
The case $K=0$ has already been excluded, so $K>0$.

Suppose $h\le0$. Then $1\le K\le R+1<P$.
If $K=1$, the identity $K=R+1+h(P-1)$ gives $P-1\mid R$,
contrary to $0<R<P-1$.
If $2\le K<P$, then
\[
 (PK-1)^2<(P^2-1)K^2+1<(PK)^2.
\]
The difference on the left is $K(2P-K)>0$, and the difference on
the right is $K^2-1>0$. Thus $a_0$ lies strictly between consecutive
squares, a contradiction. Hence $h\ge1$.
\end{proof}

It follows that the existential predicate \eqref{i17:JR} is unchanged
if $h_0$ and $i_0$ range over all integers, provided $i=i_0+1$ is
nonzero and the six coordinates $\mathbf x$ remain nonnegative.

\subsection{Seven inequalities from one square condition}

Introduce seven integer unknowns $z_1,\ldots,z_7$. To avoid confusing
the following factors with the source quantities $B_1$ and $C$, use
the symbols $\eta_j$ and $\Gamma_j$.
Define
\begin{equation}
 \begin{aligned}
 \Gamma_0&=1,\\
 d_j&=3x_j+2 &&(1\le j\le6),\\
 d_7&=3m-1,\\
 \eta_j&=1+4d_j\Gamma_{j-1}^2z_j^2,
 &\Gamma_j&=\Gamma_{j-1}\eta_j\qquad(1\le j\le7).
 \end{aligned}
 \label{i17:stack}
\end{equation}

\begin{lemma}[A coprime Pell stack]
\label{i17:pell-stack}
Subject to $z_1\cdots z_7\ne0$, the condition
$\Gamma_7\in\operatorname{Sq}(\Z)$ implies
$x_1,\ldots,x_6\ge0$ and $m>0$.
Conversely, if these seven inequalities hold, there are nonzero
$z_1,\ldots,z_7\in\Z$ for which $\Gamma_7$ is a square.
\end{lemma}
\begin{proof}
Every $\eta_j$ is $1$ modulo $4$, hence nonzero, and
$\gcd(\eta_j,\Gamma_{j-1})=1$. Thus the seven factors $\eta_j$ are
pairwise coprime. If their product is a square, unique factorization
implies that each $|\eta_j|$ is a square. The negative of an integer
square is $0$ or $3$ modulo $4$, so $\eta_j\equiv1\pmod4$ forces
each $\eta_j$ to be a positive square.

Both $\Gamma_{j-1}$ and $z_j$ are nonzero. If $d_j<0$, then
$\eta_j\le1-4=-3$, a contradiction. Hence $d_j\ge0$.
Since $d_j\equiv2\pmod3$, equality is impossible. Therefore
$d_j>0$, which is exactly $x_j\ge0$ for $1\le j\le6$ and
$m\ge1$ for $j=7$.

Conversely, every $d_j$ is then positive and congruent to $2$ modulo
$3$, so it is not a square. The Pell equation
$u_j^2-d_jv_j^2=1$ has a nontrivial positive solution.
It has one whose $v_j$ is divisible by any prescribed positive integer
$M_0$: powers of the associated unit matrix
\[
 \begin{pmatrix}u_j&d_jv_j\\v_j&u_j\end{pmatrix}
 \in\mathrm{SL}_2(\Z)
\]
eventually become the identity modulo $M_0$, because
$\mathrm{SL}_2(\Z/M_0\Z)$ is finite. A positive power remains
nontrivial and has positive lower-left entry. Set
$M_0=2|\Gamma_{j-1}|$ and choose such a solution. Then
\[
 z_j=\frac{v_j}{2\Gamma_{j-1}}\in\Z\setminus\{0\},
 \qquad \eta_j=u_j^2.
\]
Induction gives seven positive-square factors and a square product.
\end{proof}

The nonzero conditions on the seven witnesses will be imposed together
by one divisibility relation. In particular, no zero witness is being
excluded without an algebraic constraint.

\subsection{Divisibility and nonvanishing through a degree-five polynomial}

Let
\[
 \Pi(t)=(t^3-t-1)(t^2+23).
\]

\begin{lemma}[Intersectivity]
\label{i17:intersective}
The polynomial $\Pi$ has no rational root, but for every nonzero
integer $n$ there is an integer $t$ such that $n\mid\Pi(t)$.
\end{lemma}
\begin{proof}
The monic cubic $f(t)=t^3-t-1$ has no rational root: the only candidates
are $1$ and $-1$, and neither is a root. The quadratic has no real root.

It remains to construct a root modulo every prime power. The cubic has
discriminant $-23$. For an odd prime $p\ne23$, any root of the cubic
modulo $p$ is simple and therefore lifts to all powers of $p$.
If it has no root modulo $p$, it is irreducible of degree three.
Its roots lie in $\mathbb F_{p^3}$, and Frobenius permutes them by a
three-cycle. The product of their pairwise differences is fixed by this
even permutation, so lies in $\mathbb F_p$; its square is $-23$.
It is nonzero. Thus the quadratic $t^2+23$ has a simple root modulo
$p$, which also lifts to every power of $p$.
For $p=23$, the cubic has the simple root $3$, because
$f(3)=23$ and $f'(3)=26$.

For $p=2$, use $-23\equiv1\pmod8$.
If an odd $t$ satisfies $t^2\equiv-23\pmod{2^e}$ with $e\ge3$,
then one of $t$ and $t+2^{e-1}$ satisfies the same congruence modulo
$2^{e+1}$: their squares differ by $2^et$ modulo $2^{e+1}$.
The initial cases modulo $2$, $4$, and $8$ are immediate. This supplies
roots modulo every power of $2$. The Chinese remainder theorem now
gives a root modulo $|n|$.
\end{proof}

For integer arguments $d,t,v$, define
\begin{equation}
 \begin{aligned}
 \mathcal R(d,t,v)={}&
 \bigl((t+dv)^3-d^2(t+dv)-d^3\bigr)\\
 &\hspace{1.5em}\cdot\bigl((t+dv)^2+23d^2\bigr).
 \end{aligned}
 \label{i17:homogeneous}
\end{equation}
This is homogeneous of degree five in $d$ and $t$, with $v$ treated
as a parameter.

\begin{lemma}[Combining divisibility with nonvanishing]
\label{i17:merge}
For integers $d,t,z$ with $d\ne0$,
\begin{equation}
 (d\mid t\ \land\ z\ne0)
 \quad\Longleftrightarrow\quad
 \exists v\in\Z:\ zd^5\mid\mathcal R(d,t,v).
 \label{i17:merge-equivalence}
\end{equation}
Moreover, $\mathcal R(d,t,v)\ne0$ for every integer $v$.
\end{lemma}
\begin{proof}
The identity
$\mathcal R(d,t,v)=d^5\Pi(t/d+v)$ and
Lemma~\ref{i17:intersective} prove nonvanishing.
If $t=dq_0$ and $z\ne0$, choose $v$ such that
$z\mid\Pi(q_0+v)$, again using Lemma~\ref{i17:intersective}.
Then $zd^5\mid\mathcal R(d,t,v)$.

Conversely, if $z=0$, the asserted divisibility would mean
$\mathcal R(d,t,v)=0$, which is impossible.
If $d\nmid t$, choose a prime $p$ with
$\alpha=v_p(d)>\beta=v_p(t)$; here $t\ne0$.
For $r=t+dv$ we have $v_p(r)=\beta$.
In $r^3-d^2r-d^3$, the first term has the unique least valuation
$3\beta$. In $r^2+23d^2$, the first term has the unique least
valuation $2\beta$. Hence
\[
 v_p\bigl(\mathcal R(d,t,v)\bigr)
 =5\beta<5\alpha=v_p(d^5),
\]
contrary to the asserted divisibility.
\end{proof}

The general intersectivity mechanism is classical. The argument above
supplies the degree-five instance needed here without making a novelty
claim for $\Pi$ itself. Sun's recent construction uses a related
degree-six divisibility/nonvanishing device \cite[Section 5]{sun2026}.

\subsection{The actual polynomial and its square finalizer}

The quantified coordinates are exactly
\[
 \boldsymbol\xi=
 (\varepsilon,g_0,h_0,s_0,w_0,\phi,i_0,j_0,
   z_1,\ldots,z_7,v,u)\in\Z^{17}.
\]
There are eight base coordinates, seven Pell witnesses, one residue
coordinate $v$, and one final norm coordinate $u$.
Using \eqref{i17:stack}, define
\begin{align*}
 C_\ast&=\Gamma_7,& w_1&=1+C_\ast^2,\\
 w_2&=w_1(P^2+K^2+2),& Z&=(i_0+1)z_1\cdots z_7,\\
 S&=ZD_0^5,& T&=\mathcal R(D_0,T_0,v),\\
 X&=Su-T.
\end{align*}
Then put
\begin{align*}
 A_0&=S^2C_\ast,&
 A_1&=S^2w_1^2a_0,&
 A_2&=S^2w_2^2a_1,\\
 V&=X^2+A_2-A_0-A_1,&
 W&=V^2+4X^2A_2-4A_0A_1.
\end{align*}
The required fixed polynomial is
\begin{equation}
 \boxed{P_{17}=W^2-16V^2X^2A_2.}
 \label{i17:polynomial}
\end{equation}
Every intermediate symbol is a polynomial abbreviation, so none adds
a quantified coordinate. The circuit in
\path{signed17/generated/signed17.json} expands these abbreviations into
a finite division-free directed acyclic graph.

\begin{lemma}[The three-root finalizer]
\label{i17:finalizer}
With the other sixteen coordinates fixed, the assertion
\[
 \exists u\in\Z:\ P_{17}(\boldsymbol\xi)=0
\]
is equivalent to
\begin{equation}
 S\ne0,\qquad S\mid T,\qquad
 C_\ast,a_0,a_1\in\operatorname{Sq}(\Z).
 \label{i17:finalizer-conditions}
\end{equation}
In particular, these conditions are necessary for every zero; conversely,
given these conditions, an integer $u$ can be chosen.
\end{lemma}
\begin{proof}
Lemmas~\ref{i17:denominator} and~\ref{i17:merge} give $T\ne0$.
If $S=0$, then $A_0=A_1=A_2=0$, $X=-T$, and
$P_{17}=T^8\ne0$. Thus a zero must have $S\ne0$.

Choose complex roots $r_0^2=C_\ast$, $r_1^2=a_0$, and $r_2^2=a_1$.
Direct multiplication gives the polynomial identity
\begin{equation}
 P_{17}=\prod_{\epsilon\in\{-1,1\}^3}
 \bigl[X-S(\epsilon_0r_0+\epsilon_1w_1r_1+\epsilon_2w_2r_2)\bigr].
 \label{i17:norm-identity}
\end{equation}
One can verify it first by pairing the signs of $r_0,r_1$, obtaining
$(Y^2-A_0-A_1)^2-4A_0A_1$, and then substituting
$Y=X\pm Sw_2r_2$ and multiplying the two results.
Therefore a zero makes one of the signed sums in
\eqref{i17:norm-identity} equal to the rational number $u-T/S$.

For every integer $C_\ast$,
$\sqrt{|C_\ast|}\le C_\ast^2$, including $C_\ast=0$.
For all real $P,K$,
\begin{equation}
 \sqrt{\left|(P^2-1)K^2+1\right|}\le P^2+K^2+1.
 \label{i17:root-bound}
\end{equation}
Indeed,
$|(P^2-1)K^2+1|\le(P^2+1)K^2+1$, and subtraction from
$(P^2+K^2+1)^2$ leaves
$P^4+K^4+P^2K^2+2P^2+K^2\ge0$.
Consequently
\begin{equation}
 w_1>\sqrt{|C_\ast|},\qquad
 w_2>\sqrt{|C_\ast|}+w_1\sqrt{|a_0|}.
 \label{i17:separation}
\end{equation}

Consider two sign patterns that differ on at least one nonzero root,
and take the largest index at which they differ on a nonzero root.
The absolute value of that root is at least $1$, since its nonzero
square is an integer. By \eqref{i17:separation}, its weighted
contribution to the difference strictly exceeds the sum of the absolute
values of all earlier contributions. The two sums are therefore
distinct. This triangle-inequality argument applies equally to negative
radicands, whose chosen roots are complex.

The field generated by the three roots is a finite Galois extension of
$\mathbb Q$. Every automorphism changes root signs. A rational encoded
sum is fixed by every automorphism, so the separation just proved
forces each root in that sum to be fixed individually.
Every root is therefore rational. A rational number whose square is an
integer is an integer; hence $C_\ast,a_0,a_1$ are nonnegative integer
squares. Zero radicands cause no difficulty, because their roots are
already zero and their signs can be ignored.
The encoded sum is now an integer, so $T/S$ is an integer as well.
This proves necessity of \eqref{i17:finalizer-conditions}.

Conversely, choose integer roots of the three squares and set
\[
 u=T/S+r_0+w_1r_1+w_2r_2\in\Z.
\]
The corresponding factor in \eqref{i17:norm-identity} vanishes.
\end{proof}

\subsection{Equivalence to the fixed consistency statement}

\begin{theorem}[The seventeen-integer equation]
\label{i17:main}
With the coefficients fixed by the recipe above,
\[
 \mathrm{ZFC}\vdash
 \left[
 \left(\nexists\boldsymbol\xi\in\Z^{17}:P_{17}(\boldsymbol\xi)=0\right)
 \longleftrightarrow\Con(\mathrm{ZFC})
 \right].
\]
\end{theorem}
\begin{proof}
A zero gives \eqref{i17:finalizer-conditions}.
Since $D_0\ne0$, the condition $S\ne0$ gives $i\ne0$ and
$z_j\ne0$ for every $j$.
Lemma~\ref{i17:merge} gives $D_0\mid T_0$.
Lemma~\ref{i17:pell-stack} gives the six nonnegative base coordinates
and $m>0$. Lemma~\ref{i17:h-guard} then gives $h_0\ge0$.
Finally, Lemma~\ref{i17:signed-i} permits replacement of $i$ by $|i|$,
yielding $i_0\ge0$ without changing the ratio conditions.
Thus the original eight-natural-variable predicate
\eqref{i17:JR} has a solution. Its imported equivalence supplies a
finite ZFC refutation.

Conversely, a finite ZFC refutation supplies a source solution of
\eqref{i17:JR}, for which $i>0$.
Lemma~\ref{i17:pell-stack} supplies nonzero $z_1,\ldots,z_7$ and
a square $C_\ast$. Lemma~\ref{i17:merge} supplies $v$ with $S\mid T$;
also $S\ne0$. Lemma~\ref{i17:finalizer} supplies $u$, giving an
integer zero of $P_{17}$.

All new steps are elementary number theory and finite Galois theory
formalizable in ZFC. Negating the proved equivalence between existence
of a zero and existence of a refutation gives the assertion.
\end{proof}

This theorem proves the required equivalence. It does not prove either
consistency or inconsistency of ZFC.

\subsection{Exact complexity and reproducible evidence}

The delivered circuit has exactly $20\,607$ binary operations:
$672$ additions, $26$ subtractions, and $19\,909$ multiplications.
Each multiplication by a numeral and each squaring is charged.
As in package 18, construction of the fixed coefficients is outside
the arithmetic evaluation count. The $348$ named fixed leaves are all
bound by the included coefficient recipe.

The exact total degree is
\begin{equation}
 \deg P_{17}=424296\,5^{33}+53824
 =49\,394\,555\,389\,881\,134\,033\,203\,178\,824.
 \label{i17:degree}
\end{equation}
Thus this variable-count improvement has an enormous degree; the
moderate-degree $26$-integer forms retain a substantial advantage on
degree and coefficient complexity.

For a degree audit, recall $L=5^{33}$. The source degrees are
\begin{align*}
 \deg m&=2221L+11,&\deg a_0&=3800L+14,\\
 \deg a_1&=24640L+68,&\deg D_0&=6875L+20.
\end{align*}
The first six guard products have degree recurrence
$c_j=3c_{j-1}+3$, with $c_0=0$, so
$\deg\Gamma_6=1092$. The remaining degrees are
\begin{center}
\begin{tabular}{@{}lr@{}}
\toprule
Expression & Exact degree\\
\midrule
$C_\ast$ & $2221L+3289$\\
$w_1$ & $4442L+6578$\\
$w_2$ & $6342L+6586$\\
$S$ & $34375L+108$\\
$T$ & $34375L+105$\\
$X$ & $34375L+109$\\
$A_2$ & $106074L+13456$\\
\bottomrule
\end{tabular}
\end{center}
The degree of $A_2$ strictly exceeds the degrees of $X^2,A_0,A_1$.
Consequently the leading homogeneous parts of $V$, $W$, and $P_{17}$
are respectively those of $A_2$, $A_2^2$, and $A_2^4$.
The source and stack leading coefficients are nonzero, giving
\eqref{i17:degree}. The generator also propagates structural degree
bounds and a diagonal leading coefficient modulo $1\,000\,003$.
The output coefficient is $40\,718$, is nonzero, and is certified
independent of every fixed coefficient port.

The deterministic finite diagnostics include $24$ modular comparisons
using the actual original T4 coefficients, $150$ independent evaluations
of the eight-factor norm identity, roots of $\Pi$ modulo every modulus
from $1$ through $1200$, $10\,098$ consecutive-square-gap checks, and
the signed-$i$ symmetry. These checks supplement the universal proofs
above. Run
\begin{lstlisting}
python3 -B signed17/build_signed17.py
\end{lstlisting}
from the archive root. The standard-library generator writes
\path{signed17/generated/signed17.json} and
\path{signed17/generated/signed17_checks.json}.
The latter is a diagnostic receipt, not a proof-assistant certificate.

\subsection{Provenance and remaining opportunities}

The unchanged source expressions and coefficient modules come from
package 18 at revision \pin. That package records the earlier source
acquisition revision
\texttt{f0ae4e0991ca44a7bb4e5671d42e33d775e48035}.
All nineteen imported files are authenticated byte for byte against
their pinned Git blob identifiers, with SHA-256 hashes recorded in
\path{signed17/SOURCE_PINS.json}. The literal quartic and ordered tile
list also match the hashes in the upstream coefficient binding.
The generator checks every source pin before rebuilding the polynomial.
The archive preserves the upstream license and attribution files.

Several questions remain open for further simplification. Recovering
the $j_0$ or $\phi$ nonnegativity condition from the other source
conditions would remove another guard witness. A complete fixed
instantiation of the published ten-integer construction would improve
the variable count further. Optimizing the digit-polynomial circuit
could reduce its roughly $20\,500$ inherited operations. Finally, a
rebinding to the smaller $943$-tile tuple would require proof of the
complete quartic/compiler interface for that tuple; individual
coefficient ports cannot be spliced between the two bindings.
The present section makes no exhaustive novelty or absolute-minimality
claim for its general devices.


\section{Three fewer integer variables in the low-degree forms}
\label{sec:lowdegree}

Package 20 supplies exact polynomial lifts over unrestricted integers with
$54$, $45$, and $39$ variables at degrees $4$, $6$, and $8$, respectively.
Package 23 proves an independent odd-scale theorem for the signed arithmetic
interface. Their composition removes three positivity parameters in each
lift \cite{root12,separating12}.

In the source lifts, the retained scale $s$ is forced positive by a separate
residual
\[
 s-1-s_1^2-s_2^2-s_3(s_3+1).
\]
Delete this residual and the three otherwise unused parameters. Replace
$s$ everywhere else by $2v_s+1$, with $v_s\in\Z$. All remaining residuals
are unchanged except for this affine substitution. The resulting systems
have the following exact sizes:
\begin{center}
\begin{tabular}{@{}lrrrr@{}}\toprule
Form&Integer variables&Residuals&Maximum residual degree&SOS degree\\\midrule
$Q_{51}$&51&30&2&4\\
$Q_{42}$&42&21&3&6\\
$Q_{36}$&36&15&4&8\\\bottomrule
\end{tabular}
\end{center}
The full expanded residual lists are the three JSON files under
\texttt{low\_degree/generated/}. Defining $Q_n$ as the sum of their
squares completely specifies an ordinary integer polynomial.

\subsection{Why oddness is sufficient}

The imported odd-scale theorem guards only
$J,F,\alpha,z,\eta,\zeta>0$, uses an odd strong multiplier and a
square power scale, and allows the remaining coordinates to be signed.
It proves that an odd nonzero $s$ must in fact be positive at an arithmetic
solution. Its negative-$s$ branch recovers the main Pell index and obtains
one of two incompatible congruences, using explicit bounds for the
negative-parameter Pell sequence. This is a semantic theorem; simply
deleting a positivity guard would not prove it.

For completeness the compact $943$-tile tuple retains independent sufficiently
large five-power tableau widths and heights. The canonical dyadic exponent
$t=d_\ast N$ is odd, $\pc(R)=3t+2$, and the half-binomial valuation is
$v_2(Y)=\pc(R)-2=3t$. Since $Y=sq^3$ and $q=2^t$, the canonical scale
$s$ is odd. Thus every accepting tableau has the required new coordinate
$v_s=(s-1)/2$. The exact population and padding hypotheses follow from
the whole-tuple Wang adapter, not from splicing individual coefficients.

\subsection{The exposed quotients in the quartic}

The quartic lift retains quotient variables $o,j$ in place of the single
eliminated quotient $T$. Its equations give
$V=of-c=jc-R$, hence
\[
 V\equiv-c\pmod{|f|},\qquad V\equiv-R\pmod{|c|}.
\]
These are precisely the congruences used by the negative-scale proof.
Consequently that proof excludes $s<0$ before any division to recover
$T$. Once $s>0$, the ordinary strong-rank lemma gives $f^2\equiv1\pmod c$.
Since $of+R=(j+1)c$, the identity
\[
 o+Rf=f(of+R)-o(f^2-1)
\]
proves $c\mid o+Rf$. Thus $T=(o+Rf)/c$ is an integer and
$c(Tf-1)-Rf^2=of-c=V$. The full signed interface is recovered with no
unchecked divisibility assumption. The sextic and octic eliminations are
already exact integer identities and require no extra quotient argument.

The full detailed composition proof is
\texttt{low\_degree/ODD\_SCALE\_LIFTS\_PROOF.md}. Over $\Z$ a sum of
squares vanishes exactly when every residual does, proving
\[
 \mathrm{ZFC}\vdash
 (\nexists\mathbf x\in\Z^n:Q_n(\mathbf x)=0)
 \longleftrightarrow\Con(\mathrm{ZFC})
 \qquad(n=51,42,36).
\]

\subsection{Exactness of the degrees}

An affine substitution cannot increase degree. To show that degree does
not fall, the exact symbolic expansion gives
\[
 [\tau^4]Q_{51}=1,\qquad
 [J^2L_\mathrm{reg}^2v_s^2]Q_{42}=4\beta^2,\qquad
 [J^6v_s^2]Q_{36}=4\beta^6.
\]
Here $L_\mathrm{reg}$ denotes the retained register called \texttt{L} in
the sextic JSON, not the norm weight from Section~\ref{sec:definition}.
Since $\beta>0$, the degrees are exactly $4,6,8$. The sparse symbolic
sum-of-squares term counts are $342,467,845$, respectively; these are not
claims about optimally minimized circuits or numerical coefficient height.

\section{Reproducibility, proof status, and integration}
\label{sec:reproduce}

The archive is a self-contained research contribution with the article,
complete source, emitted circuits, literal coefficient data, provenance,
and verification receipts. Run
\begin{lstlisting}
python3 reproduce.py
\end{lstlisting}
from its top directory. The default run uses only Python's standard library;
LaTeX is needed only to rebuild this PDF. Optional full coefficient
materialization is controlled by the driver. All inputs required by the
default checks are included, addressing the pinned catalogue's warning
that earlier verification drivers could fail after copied files were
omitted during repository placement.

The performed checks separate three types of evidence. Exact symbolic and
integer checks verify the new polynomial definitions: independent conjugate
products, the transport shift, parity, liveness, and operation counts.
Uniform leading-term proofs establish exact degree; modular leading-form
fixtures are additional diagnostics. Byte-level checks verify the literal
coefficient payloads against their immutable manifest. None of these tests
is substituted for the semantic Pell/compiler proofs.

In particular no run purports to find an accepting K215 computation or
to decide the existence of a solution. Exhaustive bounded searches could
not prove the required absence of solutions. The new proofs are ordinary
mathematical arguments, with the precise formalization boundary stated in
Section~\ref{sec:coefficients}.

The intended repository destination is a new report directory in the
\path{con-zfc-diophantine-forms} collection. No source
repository branch, commit, or pull request was changed. The delivery keeps
the inherited sources separate from the new circuits and records one
authoritative coefficient binding for each family.

\section{Further research questions}
\label{sec:future}

\begin{enumerate}
\item \textbf{Can the index shift be optimized together with the norm?}
The useful invariant is $ac>R\Delta$, not the particular formula $q+h$.
An alternative witness parameterization that guarantees this inequality
with fewer arithmetic gates might lower the operation count as well as
degree. It must retain a proved completeness map.
\item \textbf{Can one remove another supplied root coordinate?}
The present result keeps twelve positive inputs. Eliminating an additional
coordinate requires a new integrality or congruence argument; a formal
resultant alone can introduce spurious rational roots. Any proposed
eleven-variable form should state the root-recovery order explicitly.
\item \textbf{Can the $67$-operation polynomial be improved?}
The new norm does not approach that operation count. Joint optimization
of the six-factor normalized interface, including quotient substitutions
and modular unit guards, remains a distinct problem.
\item \textbf{Can more positivity guards be inferred?}
The odd-scale theorem removes three parameters from each low-degree lift.
Identifying another coordinate whose sign is forced by the existing
equations would directly improve the quartic, sextic, and octic counts.
\item \textbf{Can the concrete signed construction approach the general
integer-variable bounds?}
The general representation theorem in Sun's September 2026 preprint uses
ten integer unknowns for arbitrary recursively enumerable sets
\cite{sun2026}. A complete fixed Con(ZFC) specialization should materialize
or finitely specify every index and coefficient, state the effective
compiler used, and compare its actual degree and description size.
\item \textbf{Which parts should be formalized first?}
The rational auxiliary-Pell lemma, the two square-comparison identities,
the noncircular bootstrap, and the explicit witness shifts are small
reusable targets. Proving a generic finite norm evaluator correct would
then connect the emitted DAG to these lemmas. The existing compiler
still needs its own end-to-end formal bridge.
\item \textbf{How should minimality be stated?}
A meaningful lower bound must specify domain, total degree, allowed
coefficient encoding, and the arithmetic model. The current collection
supports a frontier of tradeoffs, not a mathematically justified global
``simplest equation.''
\end{enumerate}

\begin{thebibliography}{99}
\raggedright
\bibitem{catalogue} V. Reshetnikov, ProveIt, \emph{Con(ZFC) in Diophantine
form}, catalogue at revision \pin.
\url{https://github.com/VladimirReshetnikov/ProveIt/tree/ae3007ca4467011e0e051c84acb4a1dd13971b0e/SetTheory/Cardinals/docs/reports/hilbert-tenth-problem/con-zfc-diophantine-forms}.
Local copy: \texttt{references/COMPARISON\_AT\_PIN.md}.
\bibitem{k215} ProveIt, \texttt{KEndpoint215.lean}, same revision,
\path{Computability/BusyBeaver/Lean/BusyBeaver/ZFC/}.
Local endpoint source: \path{references/KEndpoint215.lean}.
Literal machine SHA-256:\par
\path{fdbcef0825886035c3091c8537a69535bc96f9e265f0ae45863543d40ae70090}.
\bibitem{wang943} ProveIt wave-3 package 19, \emph{Twelve positive variables
and materialized coefficients for the consistency of ZFC}. In particular
\texttt{WANG943\_COEFFICIENTS.md} and \texttt{WANG943\_COMPACT\_LAYOUT.md},
included under \texttt{coefficients/}.
\bibitem{norm12} ProveIt wave-3 package 21, \emph{A twelve-variable equation
for Con(ZFC)}. In particular \texttt{NORM12\_PROOF.md} and
\texttt{NORM12\_LOWER\_DEGREE.md}, included under \texttt{references/}.
\bibitem{blank67} ProveIt, \emph{Independent proof and source audit:
the 67-operation blank-input polynomial}. Included as
\texttt{references/FIXED\_BLANK\_67\_PROOF.md}, together with the ordinary
and signed-domain interface proofs.
\bibitem{root12} ProveIt wave-3 package 20, \emph{Twelve positive unknowns
and materialized coefficients for ZFC consistency}. Low-degree source
polynomials and their domain reduction proof are included under
\texttt{low\_degree/upstream/20-root12/}.
\bibitem{monic12} ProveIt wave-3 package 22, \emph{Twelve-variable Diophantine
forms of ZFC consistency}. Transport shift proof:
\texttt{references/RETAINED59\_PROOF.md}.
\bibitem{separating12} ProveIt wave-3 package 23, \emph{Twelve positive
unknowns for the consistency of ZFC}. Odd-scale proof:
\texttt{low\_degree/upstream/23-separating12/proofs/ODD\_SCALE\_27\_PROOF.md}.
\bibitem{jones9} ProveIt wave-3 package 18, \emph{Three explicit Diophantine
forms of ZFC consistency}. Concrete quartic, index recipe, and Jones ratio
specialization, at the same pinned revision.
\bibitem{jones1982} J. P. Jones, \emph{Universal Diophantine equation},
Journal of Symbolic Logic \textbf{47}(3) (1982), 549--571.
\url{https://doi.org/10.2307/2273588}.
\bibitem{sun2026} Z.-W. Sun, \emph{Ten unknowns for Hilbert's tenth problem
over the integers}, arXiv:2609.01594v2, 2 September 2026, Theorem 1.1.
\url{https://arxiv.org/abs/2609.01594}.
\end{thebibliography}
\end{document}
